\documentclass[a4paper,11pt]{article}
\usepackage[utf8]{inputenc}
\usepackage[T1]{fontenc}
\usepackage[french]{babel}
\usepackage{amsmath,amssymb}
\usepackage{geometry}
\usepackage{tikz}
\usepackage{enumitem}
\usepackage{fancyhdr}
\usepackage{tcolorbox}
\usepackage{pgfplots}
\pgfplotsset{compat=1.18}

\geometry{margin=2cm}
\setlength{\headheight}{14pt}

\pagestyle{fancy}
\fancyhf{}
\lhead{BTS -- Mathématiques}
\rhead{Devoir Maison -- Séries de Fourier}
\cfoot{\thepage}

\tcbset{
  exobox/.style={colback=orange!5!white, colframe=orange!60!black, fonttitle=\bfseries, title=#1},
  reponsebox/.style={colback=gray!5!white, colframe=gray!40!black, sharp corners, boxrule=0.5pt,
                     left=4pt, right=4pt, top=4pt, bottom=4pt}
}

\newcommand{\dt}{\mathrm{d}t}

\begin{document}

\begin{center}
  {\LARGE \textbf{Devoir Maison -- Séries de Fourier}} \\[0.2cm]
  {\Large Analyse harmonique de signaux électroniques} \\[0.2cm]
  \rule{0.6\textwidth}{0.4pt} \\[0.2cm]
  {\normalsize À rendre le : \dotfill \quad Calculatrice autorisée} \\[0.1cm]
  {\normalsize Le barème est donné à titre indicatif.}
\end{center}

\vspace{0.2cm}

\noindent\textbf{Nom :} \dotfill \quad \textbf{Prénom :} \dotfill \quad \textbf{Classe :} \dotfill

\vspace{0.3cm}

\noindent\fbox{\parbox{\dimexpr\textwidth-2\fboxsep-2\fboxrule}{%
\textbf{Rappels :} Soit $f$ une fonction périodique de période $T$ et de pulsation $\omega = \dfrac{2\pi}{T}$.
Son développement en série de Fourier est :
\[
  S(t) = a_0 + \sum_{n=1}^{+\infty}\!\bigl(a_n\cos(n\omega t)+b_n\sin(n\omega t)\bigr)
\]
avec \quad
$\displaystyle a_0 = \frac{1}{T}\int_0^{T} f(t)\,\dt$,\qquad
$\displaystyle a_n = \frac{2}{T}\int_0^{T} f(t)\cos(n\omega t)\,\dt$,\qquad
$\displaystyle b_n = \frac{2}{T}\int_0^{T} f(t)\sin(n\omega t)\,\dt$.
}}

\vspace{0.4cm}

%% ================================================================
%% EXERCICE 1 — 10 points
%% ================================================================
\begin{tcolorbox}[exobox={Exercice 1 \hfill \normalfont\textit{(10 points)}}]
  \textit{Contexte :} Un générateur basse fréquence (GBF) délivre une \textbf{tension carrée symétrique}
  $v(t)$, de période $T = 2\,\mathrm{ms}$ (fréquence $f_0 = 500\,\mathrm{Hz}$) et d'amplitude $5\,\mathrm{V}$.
  On travaille en temps réduit ($t$ exprimé en ms). La tension est modélisée par :
  \[
    v(t) = \begin{cases} \phantom{-}5 & \text{si } 0 \leq t < 1 \\ -5 & \text{si } 1 \leq t < 2 \end{cases}
  \]
\end{tcolorbox}

\begin{center}
  \begin{tikzpicture}
    \begin{axis}[
      axis lines=middle,
      xlabel={$t$ (ms)},
      ylabel={$v(t)$ (V)},
      xmin=-2.5, xmax=5.5,
      ymin=-7.5, ymax=7.5,
      xtick={-2,-1,0,1,2,3,4,5},
      ytick={-5,0,5},
      width=13cm, height=5.5cm,
      grid=major,
      grid style={dashed, gray!30},
    ]
    \addplot[blue, ultra thick, const plot, mark=none]
      coordinates {(-2,5)(-1,-5)(0,5)(1,-5)(2,5)(3,-5)(4,5)(5.5,5)};
    \end{axis}
  \end{tikzpicture}
\end{center}

\begin{enumerate}[label=\textbf{\arabic*.}]
  \item Donner la valeur de la pulsation $\omega$. \hfill \textit{(1 pt)}
  \begin{tcolorbox}[reponsebox, height=1.3cm]\end{tcolorbox}

  \item Calculer $a_0$. Que représente $a_0$ pour un signal électrique ? \hfill \textit{(2 pts)}
  \begin{tcolorbox}[reponsebox, height=2cm]\end{tcolorbox}
  
  \newpage

  \item Calculer $a_n$ pour tout entier $n \geq 1$. Que peut-on conclure sur la nature de ce signal ? \hfill \textit{(2 pts)}
  \begin{tcolorbox}[reponsebox, height=4cm]\end{tcolorbox}

  \item Calculer $b_n$ pour tout entier $n \geq 1$.
        En déduire les valeurs exactes de $b_1$, $b_2$ et $b_3$. \hfill \textit{(4 pts)}
  \begin{tcolorbox}[reponsebox, height=6cm]\end{tcolorbox}

  \item Écrire le développement en série de Fourier $S(t)$ jusqu'à $n = 3$. \hfill \textit{(1 pt)}
  \begin{tcolorbox}[reponsebox, height=2cm]\end{tcolorbox}
\end{enumerate}



%% ================================================================
%% EXERCICE 2 — 10 points
%% ================================================================
\begin{tcolorbox}[exobox={Exercice 2 \hfill \normalfont\textit{(10 points)}}]
  \textit{Contexte :} Dans un \textbf{hacheur série} (convertisseur DC/DC abaisseur), la tension
  appliquée à un moteur à courant continu est un signal rectangulaire de type MLI (Modulation de
  Largeur d'Impulsion). Ce signal $u(t)$ est défini (en V, $t$ en ms) par :
  \[
    u(t) = \begin{cases} 12 & \text{si } 0 \leq t < 3 \\ \phantom{1}0 & \text{si } 3 \leq t < 4 \end{cases}
    \qquad \text{de période } T = 4\,\mathrm{ms}
  \]
  Le \textbf{rapport cyclique} vaut $\alpha = \dfrac{3}{4} = 75\,\%$ et la \textbf{tension de crête} est $E = 12\,\mathrm{V}$.
\end{tcolorbox}

\begin{enumerate}[label=\textbf{\arabic*.}]
  \item Tracer le signal $u$ sur $[-4\,;\,8]$ dans le repère ci-dessous. \hfill \textit{(1,5 pts)}

  \begin{center}
    \begin{tikzpicture}
      \begin{axis}[
        axis lines=middle,
        xlabel={$t$ (ms)},
        ylabel={$u(t)$ (V)},
        xmin=-4.5, xmax=8.5,
        ymin=-2, ymax=14.5,
        xtick={-4,-3,-2,-1,0,1,2,3,4,5,6,7,8},
        ytick={0,4,8,12},
        width=14cm, height=5cm,
        grid=major,
        grid style={dashed, gray!30},
      ]
      \end{axis}
    \end{tikzpicture}
  \end{center}

  \item Donner la valeur de la pulsation $\omega$. \hfill \textit{(0,5 pt)}
  \begin{tcolorbox}[reponsebox, height=1.2cm]\end{tcolorbox}

  \item Calculer $a_0$. Vérifier que $a_0 = \alpha E$ et interpréter ce résultat en termes de tension électrique. \hfill \textit{(2 pts)}
  \begin{tcolorbox}[reponsebox, height=2.5cm]\end{tcolorbox}

  \item Calculer $a_n$ pour tout entier $n \geq 1$.
        En déduire les valeurs exactes de $a_1$, $a_2$ et $a_3$. \hfill \textit{(3 pts)}
  \begin{tcolorbox}[reponsebox, height=6cm]\end{tcolorbox}

  \item Calculer $b_n$ pour tout entier $n \geq 1$.
        En déduire les valeurs exactes de $b_1$, $b_2$ et $b_3$. \hfill \textit{(3 pts)}
  \begin{tcolorbox}[reponsebox, height=6cm]\end{tcolorbox}

\end{enumerate}

\vfill
\begin{center}
  \fbox{\textbf{Tourner la page $\longrightarrow$}}
\end{center}

\newpage

%% ================================================================
%% BONUS — Illustration graphique et interprétation
%% ================================================================
\begin{tcolorbox}[exobox={Bonus -- Illustration graphique et interprétation}]
  On donne ci-dessous le tracé de la somme partielle $S_5(t)$, obtenue en conservant les cinq premiers
  harmoniques non nuls de la série calculée à l'exercice 1 :
  \[
    S_5(t) = \frac{20}{\pi}\sin(\pi t) + \frac{20}{3\pi}\sin(3\pi t) + \frac{20}{5\pi}\sin(5\pi t)
  \]
  Le signal original $v(t)$ est tracé en pointillé rouge pour référence.
\end{tcolorbox}

\vspace{0.3cm}

\begin{center}
  \begin{tikzpicture}
    \begin{axis}[
      axis lines=middle,
      xlabel={$t$ (ms)},
      ylabel={},
      xmin=-2.3, xmax=5.3,
      ymin=-7.8, ymax=7.8,
      xtick={-2,-1,0,1,2,3,4,5},
      ytick={-5,0,5},
      width=13cm, height=5.5cm,
      grid=major,
      grid style={dashed, gray!30},
      legend style={at={(1.02,1)}, anchor=north west},
    ]
    %% Signal carré original (pointillé rouge)
    \addplot[red, dashed, thick, mark=none] coordinates {
      (-2,5)(-1,5)(-1,-5)(0,-5)(0,5)(1,5)(1,-5)(2,-5)(2,5)(3,5)(3,-5)(4,-5)(4,5)(5.3,5)
    };
    %% Approximation S5 (bleu)
    \addplot[blue, thick, domain=-2.3:5.3, samples=600]
      { 20/pi * sin(deg(pi*x)) + 20/(3*pi) * sin(deg(3*pi*x)) + 20/(5*pi) * sin(deg(5*pi*x)) };
    \legend{$v(t)$ (signal original), $S_5(t)$ (approximation)}
    \end{axis}
  \end{tikzpicture}
\end{center}

\vspace{0.2cm}

\begin{enumerate}[label=\textbf{\arabic*.}]
  \item Commenter l'allure de la courbe $S_5(t)$ : vers quelle fonction $v$ semble-t-elle converger
        au fur et à mesure que l'on ajoute des harmoniques ? Que se passe-t-il au voisinage des
        discontinuités (phénomène observé) ?
  \begin{tcolorbox}[reponsebox, height=3cm]\end{tcolorbox}

  \item \textbf{Application au filtrage (signal MLI de l'exercice 2) :}

  Un ingénieur branche le signal $u(t)$ de l'exercice 2 à l'entrée d'un filtre passe-bas
  de fréquence de coupure $f_c = 20\,\mathrm{Hz}$.

  \begin{enumerate}[label=\textbf{\alph*.}]
    \item Quelle est la fréquence fondamentale $f_1$ du signal $u(t)$ ? \hfill \textit{(0,5 pt)}
    \begin{tcolorbox}[reponsebox, height=1.2cm]\end{tcolorbox}

    \item Justifier que toutes les composantes harmoniques ($f_1$, $2f_1$, $3f_1$, \ldots) sont
          fortement atténuées par ce filtre. \hfill \textit{(1 pt)}
    \begin{tcolorbox}[reponsebox, height=2cm]\end{tcolorbox}

    \item Quelle tension retrouve-t-on en sortie du filtre ? Quel est l'intérêt de ce résultat
          dans le contexte du hacheur ? \hfill \textit{(0,5 pt)}
    \begin{tcolorbox}[reponsebox, height=2cm]\end{tcolorbox}
  \end{enumerate}
\end{enumerate}

\end{document}
